\section{What a fixed observation cap can identify}
\label{sec:identifiability}

The exact dual also identifies an obstruction that is invisible in a
small-cap numerical upper bound: increasing the number of units cannot
resolve two priors with identical cap-$N$ one-unit distributions.

Let $L=\lceil N/2\rceil$ and define the moment-matching diameter
\begin{equation}\label{eq:cap-diameter}
 \Delta_N=\max\left\{
 \theta(q_H)-\theta(q_L):q_L,q_H\in\mathcal Q_\rho,\
 \int(q_H-q_L)d^{2j+1}\,d\mu=0\ (0\le j<L)
 \right\}.
\end{equation}
Compactness ensures that this maximum is attained.

\begin{theorem}[Exact fixed-cap bias floor]\label{thm:cap-identifiability}
For every magnitude law,
\begin{equation}\label{eq:cap-bias-floor}
 \lim_{B\to\infty}P_N(B)=\Delta_N/2.
\end{equation}
If $\mu$ has exactly $m$ distinct positive support points, then
$\Delta_N=0$ if and only if $N\ge2m-1$.
If $\mu$ has infinite support, then $\Delta_N>0$ for every finite $N$.
\end{theorem}

\begin{proof}
The distribution of $N$ Bernoulli observations from one unit determines
the first $N$ moments of its latent Bernoulli mean. The difference of the
$k$th such moments under two priors is
\[
 \int(q_H-q_L)
       \{(1/2+d)^k-(1/2-d)^k\}\,d\mu.
\]
These polynomials span exactly the odd monomials
$d,d^3,\ldots,d^{2L-1}$. Thus the constraints in
\eqref{eq:cap-diameter} are equivalent to equality of the full depth-$N$
count distributions. Equality at depth $N$ implies equality for every
stopped experiment with that cap, by marginalization and independent
randomization. Conversely, the policy that always observes exactly $N$
samples witnesses a positive information numerator whenever those count
distributions differ. Therefore $K_N=0$ exactly for pairs satisfying
\eqref{eq:cap-diameter}; otherwise $K_N>0$ or $K_N=\infty$.

Apply the decreasing compact-set argument from
Lemma~\ref{lem:infinite-dual} to the dual ratios as $B\to\infty$.
Their pointwise limit is $(\Delta\theta)_+^2/4$ when $K_N=0$, and zero
otherwise. Theorem~\ref{thm:exact-dual} gives
\eqref{eq:cap-bias-floor}.

For the support assertion, consider the finite-dimensional subspace
\[
 \mathcal V_L=\operatorname{span}\{d,d^3,\ldots,d^{2L-1}\}
 \subset L^2(\mu).
\]
If the support has $m\le L$ points, polynomial interpolation gives
$p(d)=dQ(d^2)=1$ on that support, with $\deg Q\le m-1$.
Then $1\in\mathcal V_L$, so moment matching forces zero target gap.

If the support has more than $L$ points, $1\notin\mathcal V_L$.
Indeed, a polynomial $Q$ of degree at most $L-1$ cannot agree with
$x^{-1/2}$ at $L+1$ distinct positive values: otherwise repeated Rolle's
theorem would force the nonvanishing $L$th derivative of $x^{-1/2}$ to
vanish somewhere. The same argument covers infinite support, since an
almost-everywhere polynomial identity extends to the support.

Let $h=1-\operatorname{proj}_{\mathcal V_L}1$. It is a nonzero bounded
function, orthogonal to every required odd monomial, and
\[
 \int h\,d\mu=\int h^2\,d\mu>0.
\]
Set $\epsilon=\rho/(4\|h\|_\infty)$ and
$q_H=\rho/2+\epsilon h$, $q_L=\rho/2-\epsilon h$.
Both priors lie pointwise between $\rho/4$ and $3\rho/4$ and hence
belong to $\mathcal Q_\rho$. They match the odd moments but have positive
target gap $2\epsilon\int h^2\,d\mu$. This proves strict positivity.
Finally $L\ge m$ is equivalent to $N\ge2m-1$.
\end{proof}

The proof also yields the explicit lower estimate
\[
 \Delta_N/2\ge
 \frac{\rho\|h\|_{L^2(\mu)}^2}{4\|h\|_{L^\infty(\mu)}}
 \quad\text{when }1\notin\mathcal V_L.
\]
This theorem identifies the target functional from cap-limited laws;
it does not require recovery of the entire allocation $q$.
This theorem concerns the fixed-cap program and its infinite-budget limit.
It does not impose a bias floor on the unrestricted experiment. The
monotone brackets in Section~\ref{sec:genie} increase the cap when needed,
and their stopping criterion tests the remaining gap directly.
